\chapter{Integreren}\label{ch:g12:integ}

Integreren beantwoordt twee vragen tegelijk: wat is de oppervlakte onder een
kromme, en hoe haal je een \omterm{def:g11:func:function}{functie} terug uit haar veranderingssnelheid? De
hoofdstelling van de integraalrekening zegt dat dat dezelfde vraag is --- de
diepste en nuttigste ontdekking van de zeventiende-eeuwse wiskunde.

\section{De integraal van een continue functie}

\begin{definition}[Integraal van een positieve functie]\label{def:g12:integ:area}
Zij $f$ \omterm{def:g12:limcont:continuity}{continu} en positief op $\intcc{a}{b}$. De
\emph{integraal}\index{integraal}
\[
\int_a^b f(x)\,\dd x
\]
is de oppervlakte, in oppervlakte-eenheden, van het gebied begrensd door de
kromme van $f$, de $x$-as en de verticale rechten $x = a$ en $x = b$.
\end{definition}

\begin{omfigure}
\begin{tikzpicture}
\begin{axis}[
  width=8.4cm, height=5cm,
  axis lines=middle, xmin=-0.4, xmax=3.4, ymin=-0.3, ymax=2.3,
  xtick={0.5,2.7}, xticklabels={$a$,$b$}, ytick=\empty,
  samples=100, domain=-0.2:3.3]
  \addplot[name path=curve, omDef, thick, domain=0.2:3.2]
    {1 + 0.5*sin(deg(2*x)) + 0.2*x};
  \path[name path=axis] (0.5,0) -- (2.7,0);
  \addplot[omDef!20] fill between[of=curve and axis, soft clip={domain=0.5:2.7}];
  \node at (axis cs:1.6,0.55) {$\displaystyle\int_a^b f$};
\end{axis}
\end{tikzpicture}
\end{omfigure}

Voor een \omterm{def:g11:func:function}{functie} met willekeurig teken tellen de oppervlakten onder de
$x$-as negatief mee; en men stelt
$\int_b^a f(x)\,\dd x = -\int_a^b f(x)\,\dd x$.

\begin{omfigure}
\begin{tikzpicture}
\begin{axis}[omaxis,
  width=9.6cm, height=4.8cm,
  xmin=-0.5, xmax=7, ymin=-1.4, ymax=1.4,
  xtick={3.1416,6.2832}, xticklabels={$\pi$,$2\pi$}, ytick={-1,1}]
  \addplot[name path=sine, omDef, thick, domain=0:6.2832, samples=120]
    {sin(deg(x))};
  \path[name path=xaxis] (0,0) -- (6.2832,0);
  \addplot[omDef!20] fill between[of=sine and xaxis,
    soft clip={domain=0:3.1416}];
  \addplot[omThm!20] fill between[of=sine and xaxis,
    soft clip={domain=3.1416:6.2832}];
  \node at (axis cs:1.57,0.4) {$+$};
  \node at (axis cs:4.71,-0.4) {$-$};
\end{axis}
\end{tikzpicture}

{\small Getekende oppervlakten: $\int_0^{2\pi} \sin x\,\dd x = 0$, want het
gebied onder de as (rood) heft dat erboven (blauw) op.}
\end{omfigure}

\begin{proposition}[Eigenschappen van de integraal]\label{prop:g12:integ:properties}
Zij $f, g$ \omterm{def:g12:limcont:continuity}{continu} op een \omterm{def:g10:numbers:interval}{interval} dat $a, b, c$ bevat, en $\lambda \in \R$.
\begin{enumerate}
  \item \emph{Lineariteit:}
        $\displaystyle\int_a^b (f + \lambda g) = \int_a^b f + \lambda \int_a^b g$.
  \item \emph{Relatie van Chasles:}
        $\displaystyle\int_a^c f = \int_a^b f + \int_b^c f$.
  \item \emph{Positiviteit:} is $a \leq b$ en $f \geq 0$ op $\intcc{a}{b}$,
        dan is $\displaystyle\int_a^b f \geq 0$; is $f \leq g$, dan is
        $\displaystyle\int_a^b f \leq \int_a^b g$.
\end{enumerate}
\end{proposition}

\begin{proof}
Chasles en de positiviteit volgen uit de duiding als oppervlakte (oppervlakten
tellen op wanneer gebieden naast elkaar liggen). De ongelijkheid volgt door de
positiviteit op $g - f$ toe te passen. De lineariteit is intuïtief duidelijk
(stapel de oppervlakten) en wordt aan de universiteit streng bewezen;
\emph{hier wordt ze aangenomen}.
\end{proof}

\section{De hoofdstelling van de integraalrekening}

\begin{definition}[Primitieve functie]\label{def:g12:integ:primitive}
Een \emph{primitieve functie}\index{primitieve functie} van $f$ op een
\omterm{def:g10:numbers:interval}{interval} $I$ is een differentieerbare \omterm{def:g11:func:function}{functie} $F$ op $I$ waarvoor $F' = f$.
\end{definition}

\begin{proposition}\label{prop:g12:integ:primunique}
Is $F$ een \omterm{def:g12:integ:primitive}{primitieve functie} van $f$ op een \omterm{def:g10:numbers:interval}{interval} $I$, dan zijn de
\omterm{def:g12:integ:primitive}{primitieve functies} van $f$ op $I$ precies de \omterm{def:g11:func:function}{functies} $F + c$ met
$c \in \R$. Bij gegeven $x_0 \in I$ en $y_0 \in \R$ is er precies één
\omterm{def:g12:integ:primitive}{primitieve functie} met $F(x_0) = y_0$.
\end{proposition}

\begin{proof}
Is $G' = F' = f$, dan is $(G - F)' = 0$ op het \omterm{def:g10:numbers:interval}{interval} $I$, dus is $G - F$
constant.\footnote{Dat een \omterm{def:g11:func:function}{functie} met \omterm{def:g12:deriv:derivative}{afgeleide} nul op een \omterm{def:g10:numbers:interval}{interval} constant
is, volgt uit \cref{thm:g12:deriv:variations}: ze is zowel \omterm{def:g12:seq:monotonic}{stijgend} als
\omterm{def:g10:functions:variations}{dalend}.} Omgekeerd is elke $F + c$ een \omterm{def:g12:integ:primitive}{primitieve functie}. De voorwaarde
$F(x_0) = y_0$ legt de constante vast.
\end{proof}

\begin{theorem}[Hoofdstelling van de integraalrekening]\label{thm:g12:integ:ftc}
\index{hoofdstelling van de integraalrekening}
Zij $f$ \omterm{def:g12:limcont:continuity}{continu} op een \omterm{def:g10:numbers:interval}{interval} $I$ en $a \in I$. De \omterm{def:g11:func:function}{functie}
\[
F \colon x \longmapsto \int_a^x f(t)\,\dd t
\]
is de \omterm{def:g12:integ:primitive}{primitieve functie} van $f$ op $I$ die in $a$ verdwijnt. Bijgevolg geldt
voor elke \omterm{def:g12:integ:primitive}{primitieve functie} $G$ van $f$ en alle $a, b \in I$:
\[
\int_a^b f(t)\,\dd t = \bigl[G(t)\bigr]_a^b = G(b) - G(a).
\]
\end{theorem}

\begin{omfigure}
\begin{tikzpicture}
\begin{axis}[omaxis,
  width=9.2cm, height=5.6cm,
  xmin=-0.4, xmax=4.2, ymin=-0.4, ymax=2.9,
  xtick={0.5,2.2,2.7}, xticklabels={$a$,$x$,\ $x+h$}, ytick=\empty]
  \addplot[name path=curve, omDef, thick, domain=0:3.9, samples=80]
    {1 + 0.4*x + 0.2*sin(deg(2*x))};
  \path[name path=xaxis] (0,0) -- (3.9,0);
  \addplot[omDef!20] fill between[of=curve and xaxis,
    soft clip={domain=0.5:2.2}];
  \addplot[omProp!40] fill between[of=curve and xaxis,
    soft clip={domain=2.2:2.7}];
  \draw[black, dashed, thin] (axis cs:2.2,1.69) -- (axis cs:2.7,1.69);
  \draw[black, dashed, thin] (axis cs:2.2,1.925) -- (axis cs:2.7,1.925)
    node[above right, font=\small] {$f(x+h)$};
  \node at (axis cs:1.35,0.6) {$F(x)$};
\end{axis}
\end{tikzpicture}

{\small Het idee van het bewijs: de aangroei $F(x+h) - F(x)$ is de oppervlakte
van de smalle oranje strook, ingeklemd tussen rechthoeken met hoogten $f(x)$
en $f(x+h)$ op basis $h$.}
\end{omfigure}

\begin{proof}[Bewijs voor stijgende $f$]
Kies $x \in I$ en $h > 0$ met $x + h \in I$. Volgens Chasles is
\[
F(x+h) - F(x) = \int_x^{x+h} f(t)\,\dd t .
\]
Omdat $f$ \omterm{def:g12:seq:monotonic}{stijgend} is, geldt $f(x) \leq f(t) \leq f(x+h)$ voor
$t \in \intcc{x}{x+h}$, en door integralen te vergelijken (de \omterm{def:g12:integ:area}{integraal} van
een constante $c$ over een \omterm{def:g10:numbers:interval}{interval} van lengte $h$ is $ch$):
\[
h\,f(x) \leq F(x+h) - F(x) \leq h\,f(x+h),
\]
dus
\[
f(x) \leq \frac{F(x+h) - F(x)}{h} \leq f(x+h).
\]
Voor $h \to 0^+$ is $f(x+h) \to f(x)$ door de \omterm{def:g12:limcont:continuity}{continuïteit}, en de
insluitstelling geeft dat het differentiequotiënt naar $f(x)$ streeft; het
geval $h < 0$ verloopt symmetrisch. Bijgevolg is $F' = f$ en $F(a) = 0$. Het
algemene (niet-monotone) continue geval wordt aan de universiteit bewezen.

Ten slotte: is $G$ een willekeurige \omterm{def:g12:integ:primitive}{primitieve functie}, dan is $G = F + c$
(\cref{prop:g12:integ:primunique}), zodat
$G(b) - G(a) = F(b) - F(a) = \int_a^b f$.
\end{proof}

\begin{example}
$\displaystyle\int_0^1 x^2\,\dd x = \left[\frac{x^3}{3}\right]_0^1 =
\frac13$: de oppervlakte onder de \omterm{def:g11:quad:parabola}{parabool} is een derde van het
eenheidsvierkant, zoals Archimedes wist.
\end{example}

De tabel van \omterm{def:g12:integ:primitive}{primitieve functies} lees je af van de tabel van \omterm{def:g12:deriv:derivative}{afgeleiden}
($u$ is een differentieerbare \omterm{def:g11:func:function}{functie}, $c$ een willekeurige constante):
\begin{center}
\renewcommand{\arraystretch}{1.7}
\begin{tabular}{c|c||c|c}
$f(x)$ & primitieve & $f$ & primitieve \\
\hline
$x^n \ (n \neq -1)$ & $\dfrac{x^{n+1}}{n+1} + c$ &
$u'u^n \ (n \neq -1)$ & $\dfrac{u^{n+1}}{n+1} + c$\\
$\dfrac1x \ (x > 0)$ & $\ln x + c$ &
$\dfrac{u'}{u} \ (u > 0)$ & $\ln u + c$\\
$\eu^x$ & $\eu^x + c$ &
$u'\eu^u$ & $\eu^u + c$\\
$\cos x$ & $\sin x + c$ &
$\dfrac{u'}{\sqrt u} \ (u>0)$ & $2\sqrt u + c$\\
$\sin x$ & $-\cos x + c$ & &
\end{tabular}
\end{center}

\begin{method}[De vorm $u' \times (\text{iets in } u)$ herkennen]\label{met:g12:integ:uprime}
Zoek bij het integreren van een product naar een factor die, op een
vermenigvuldigende constante na, de \omterm{def:g12:deriv:derivative}{afgeleide} van een binnenfunctie $u$ is.
Zo is in $\int_0^1 x\,\eu^{x^2}\dd x$ de factor $x$ gelijk aan
$\frac12 (x^2)'$:
\[
\int_0^1 x\,\eu^{x^2}\dd x = \frac12\left[\eu^{x^2}\right]_0^1
= \frac{\eu - 1}{2}.
\]
\end{method}

\begin{theorem}[Partiële integratie]\label{thm:g12:integ:ibp}
\index{partiële integratie}
Zij $u, v$ \omterm{def:g12:deriv:derivative}{differentieerbaar} op $\intcc{a}{b}$ met continue \omterm{def:g12:deriv:derivative}{afgeleiden}. Dan
geldt
\[
\int_a^b u'(t)\,v(t)\,\dd t
= \bigl[u(t)\,v(t)\bigr]_a^b - \int_a^b u(t)\,v'(t)\,\dd t .
\]
\end{theorem}

\begin{proof}
De productregel geeft $(uv)' = u'v + uv'$; beide leden over $\intcc{a}{b}$
integreren en de hoofdstelling op het linkerlid toepassen levert
$\bigl[uv\bigr]_a^b = \int_a^b u'v + \int_a^b uv'$.
\end{proof}

\begin{example}\label{ex:g12:integ:ibp}
$\displaystyle\int_0^1 t\,\eu^{t}\,\dd t$: neem $u' = \eu^t$ en $v = t$, dus
$u = \eu^t$ en $v' = 1$:
\[
\int_0^1 t\,\eu^t\,\dd t = \bigl[t\,\eu^t\bigr]_0^1 - \int_0^1 \eu^t\,\dd t
= \eu - (\eu - 1) = 1 .
\]
\end{example}

\section{Toepassingen}

\begin{definition}[Gemiddelde waarde]\label{def:g12:integ:mean}
De \emph{gemiddelde waarde}\index{gemiddelde waarde} van een continue \omterm{def:g11:func:function}{functie}
$f$ op $\intcc{a}{b}$ ($a<b$) is
\[
\mu = \frac{1}{b-a}\int_a^b f(t)\,\dd t .
\]
\end{definition}

\begin{proposition}\label{prop:g12:integ:meanbounds}
Is $m \leq f \leq M$ op $\intcc{a}{b}$, dan is $m \leq \mu \leq M$.
\end{proposition}

\begin{proof}
Integreer de ongelijkheden $m \leq f(t) \leq M$ over $\intcc{a}{b}$ en deel
door $b - a > 0$.
\end{proof}

\begin{method}[Oppervlakte tussen twee krommen]\label{met:g12:integ:areabetween}
Is $f \geq g$ op $\intcc{a}{b}$, dan is de oppervlakte tussen de twee krommen
gelijk aan $\int_a^b \bigl(f(x) - g(x)\bigr)\dd x$. Kruisen de krommen elkaar,
splits het \omterm{def:g10:numbers:interval}{interval} dan in de snijpunten en integreer $\abs{f - g}$ stuk voor
stuk.
\end{method}

\begin{omfigure}
\begin{tikzpicture}
\begin{axis}[omaxis,
  width=8.4cm, height=5.8cm,
  xmin=-2.1, xmax=2.9, ymin=-1.8, ymax=5.6,
  xtick={-1,2}, ytick=\empty]
  \addplot[name path=line, omThm, thick, domain=-1.7:2.5] {x + 1};
  \addplot[name path=parab, omDef, thick, domain=-1.9:2.45] {x^2 - 1};
  \addplot[omDef!20] fill between[of=line and parab,
    soft clip={domain=-1:2}];
  \fill (axis cs:-1,0) circle (1.5pt);
  \fill (axis cs:2,3) circle (1.5pt);
\end{axis}
\end{tikzpicture}

{\small De oppervlakte tussen de rechte $y = x + 1$ (rood) en de \omterm{def:g11:quad:parabola}{parabool}
$y = x^2 - 1$ (blauw) is $\int_{-1}^{2}
\bigl((x+1) - (x^2-1)\bigr)\dd x = \frac92$.}
\end{omfigure}

\section{Oefeningen}

\begin{exercise}[$\star$]\label{exo:g12:integ:1}
Bereken
\[
\int_1^2 \left(3x^2 - \frac{1}{x^2}\right)\dd x, \qquad
\int_0^{\pi/2} \cos t \,\dd t, \qquad
\int_0^{1} \frac{\dd t}{2t+1} .
\]
\end{exercise}

\begin{exercise}[$\star$]\label{exo:g12:integ:2}
Zoek de \omterm{def:g12:integ:primitive}{primitieve functie} $F$ van $f(x) = x\eu^{x^2}$ op $\R$ waarvoor
$F(0) = 1$.
\end{exercise}

\begin{exercise}[$\star$]\label{exo:g12:integ:3}
Bereken de \omterm{def:g12:integ:mean}{gemiddelde waarde} van $f(t) = \sin t$ op $\intcc{0}{\pi}$, en duid
het resultaat op een \omterm{def:g10:functions:graph}{grafiek}.
\end{exercise}

\begin{exercise}[$\star\star$]\label{exo:g12:integ:4}
Bereken met \omterm{thm:g12:integ:ibp}{partiële integratie}
\[
\int_1^{\eu} \ln t\,\dd t
\qquad\text{en}\qquad
\int_0^{\pi} t \sin t\,\dd t .
\]
\end{exercise}

\begin{exercise}[$\star\star$]\label{exo:g12:integ:5}
Bereken de oppervlakte van het gebied tussen de \omterm{def:g11:quad:parabola}{parabool} $y = x^2$ en de
rechte $y = x + 2$.
\end{exercise}

\begin{exercise}[$\star\star$]\label{exo:g12:integ:6}
Zij $I = \displaystyle\int_0^1 \frac{\dd t}{1 + t}$.
\begin{enumerate}
  \item Bereken $I$.
  \item Zij voor $n \in \N$ het getal
        $I_n = \displaystyle\int_0^1 \frac{t^n}{1+t}\dd t$. Toon aan dat
        $I_n + I_{n+1} = \dfrac{1}{n+1}$ en dat
        $0 \leq I_n \leq \dfrac{1}{n+1}$.
  \item Leid af dat
        $1 - \frac12 + \frac13 - \dots + \frac{(-1)^{n-1}}{n}
        \xrightarrow[n \to +\infty]{} \ln 2$.
\end{enumerate}
\end{exercise}

\begin{exercise}[$\star\star$]\label{exo:g12:integ:7}
De snelheid van een trein (in m/s) gedurende de eerste $100$ seconden na het
vertrek wordt gemodelleerd door $v(t) = 30\bigl(1 - \eu^{-t/50}\bigr)$.
Bereken de afgelegde afstand tijdens die $100$ seconden, en de \omterm{def:g11:stat:mean}{gemiddelde}
snelheid van de trein over dat \omterm{def:g10:numbers:interval}{interval}.
\end{exercise}

\begin{exercise}[$\star\star\star$]\label{exo:g12:integ:8}
Zij voor $n \geq 1$ de som
$S_n = \dfrac1n \displaystyle\sum_{k=1}^{n} \frac{1}{1 + k/n}$.
\begin{enumerate}
  \item Duid $S_n$ als een oppervlakte van rechthoeken die een gebied onder
        de kromme van $t \mapsto \frac{1}{1+t}$ op $\intcc{0}{1}$ benaderen.
  \item Toon met de monotonie van $t \mapsto \frac{1}{1+t}$ aan dat
        \[
        S_n \leq \int_0^1 \frac{\dd t}{1+t} \leq S_n + \frac1n\left(1 - \frac12\right),
        \]
        en leid $\lim\limits_{n\to+\infty} S_n = \ln 2$ af.
\end{enumerate}
\end{exercise}

\begin{exercise}[$\star\star\star$]\label{exo:g12:integ:9}
\emph{(Integralen van Wallis.)} Zij voor $n \in \N$ het getal
$W_n = \displaystyle\int_0^{\pi/2} \sin^n t\,\dd t$.
\begin{enumerate}
  \item Bereken $W_0$ en $W_1$.
  \item Toon, door $\sin^{n+2}t = \sin t \cdot \sin^{n+1} t$ te schrijven en
        partieel te integreren, aan dat $W_{n+2} = \dfrac{n+1}{n+2}\,W_n$.
  \item Leid $W_2$, $W_3$ en $W_4$ af, en toon aan dat $(W_n)$ \omterm{def:g10:functions:variations}{dalend} en
        positief is.
\end{enumerate}
\end{exercise}

\section{Opgave: Archimedes tegen de machine}

\begin{problem}[{Weekendopgave --- de oppervlakte onder de parabool, op drie
manieren berekend over tweeëntwintig eeuwen, en de geheime identiteit van de
logaritme als oppervlakte}]
\label{pb:g12:integ:1}
Rond 240 v.Chr. berekende Archimedes de exacte oppervlakte van een
parabolisch segment --- zonder \omterm{def:g10:coordgeom:system}{coördinaten}, zonder limieten, zonder algebra.
Negentien eeuwen later deden de rechthoeken van Riemann het over met brute
insluiting; en de hoofdstelling van de integraalrekening
(\cref{thm:g12:integ:ftc}) doet het vandaag in één regel. Deze opgave speelt
alle drie de wedstrijden, en gebruikt daarna dezelfde machine om te onthullen
wat de \omterm{def:g12:exp:ln}{logaritme} werkelijk is: een oppervlakte met een schaalsymmetrie.

\medskip
\noindent\textbf{Deel I --- Vlot rekenen.}
\begin{enumerate}
  \item Bereken $\displaystyle\int_0^1 (3x^2 - 2x + 1)\,\dd x$, \
        $\displaystyle\int_1^{\eu} \frac{\dd x}{x}$ \ en
        $\displaystyle\int_0^{\pi/2} \cos x\,\dd x$.
  \item Spot de vormen $u'u$ (\cref{met:g12:integ:uprime}):
        $\displaystyle\int_0^1 x\,\eu^{x^2}\dd x$ en
        $\displaystyle\int_0^1 \frac{2x}{x^2 + 1}\,\dd x$.
  \item Met \omterm{thm:g12:integ:ibp}{partiële integratie} (\cref{thm:g12:integ:ibp}):
        $\displaystyle\int_0^1 x\,\eu^x \dd x$ en
        $\displaystyle\int_1^{\eu} \ln x\,\dd x$.
  \item Bereken de \omterm{def:g12:integ:mean}{gemiddelde waarde} (\cref{def:g12:integ:mean}) van $\sin$
        over $\intcc{0}{\pi}$ --- en merk op dat ze \emph{niet} $\frac12$ is.
  \item Bereken de oppervlakte tussen de rechte $y = x$ en de \omterm{def:g11:quad:parabola}{parabool}
        $y = x^2$ over $\intcc{0}{1}$ (\cref{met:g12:integ:areabetween}).
\end{enumerate}

\medskip
\noindent\textbf{Deel II --- De \omterm{def:g11:quad:parabola}{parabool}, op drie manieren.}
\begin{enumerate}[resume]
  \item De opzet van Riemann voor de oppervlakte onder $y = x^2$ op
        $\intcc{0}{1}$: schrijf de ondersom $L_n$ en de bovensom $U_n$ over
        $n$ gelijke rechthoeken op (zoals in \cref{exo:g12:integ:8}).
  \item Bewijs met inductie (de machine van \cref{pb:g12:seq:1}) de formule
        voor de som van de kwadraten:
        \[
        1^2 + 2^2 + \dots + n^2 = \frac{n(n+1)(2n+1)}{6} .
        \]
  \item Leid gesloten uitdrukkingen voor $U_n$ en $L_n$ af, bereken hun
        gemeenschappelijke limiet, en besluit: de oppervlakte is $\frac13$.
  \item En nu de machine: bereken $\int_0^1 x^2\,\dd x$ met de
        hoofdstelling, in één regel. Vergelijk de inspanningen.
  \item Archimedes formuleerde het anders: \emph{een parabolisch segment is
        $\frac43$ van zijn ingeschreven driehoek}. Bereken voor het segment
        dat de koorde van $(-1, 1)$ naar $(1, 1)$ uit $y = x^2$ snijdt: de
        oppervlakte van het segment met een \omterm{def:g12:integ:area}{integraal}, de oppervlakte van de
        ingeschreven driehoek (met top in $(0, 0)$), en ga de verhouding van
        de meester na.
  \item De eigen methode van Archimedes: vul het segment met de grote
        driehoek $T$, daarna met twee driehoeken die samen $\frac T4$ meten,
        daarna met vier die samen $\frac{T}{16}$ meten, enzovoort. Sommeer de
        meetkundige reeks en haal $\frac43 T$ terug --- de eindeloos afgebeten
        chocoladereep uit het onderbouwvolume, opgegeten door een Griekse
        meetkundige.
  \item Telkens één zin: uitputting (Archimedes), insluiting (Riemann),
        terugafleiden (Newton en Leibniz) --- wat heeft elk nodig, en wat
        geeft elk?
\end{enumerate}

\medskip
\noindent\textbf{Deel III --- De \omterm{def:g12:exp:ln}{logaritme} is een oppervlakte.}
Stel voor $x > 0$: $A(x) = \displaystyle\int_1^x \frac{\dd t}{t}$.
\begin{enumerate}[resume]
  \item Geef $A(1)$ en $A'(x)$ (\cref{thm:g12:integ:ftc}), en besluit dat $A$
        precies de \omterm{def:g12:exp:ln}{natuurlijke logaritme} van \cref{def:g12:exp:ln} is.
  \item Het schaalwonder: kies $a > 0$ vast en bestudeer
        $g(x) = A(ax) - A(x)$. Bereken $g'$ (kettingregel), leid af dat $g$
        constant is, bepaal die constante --- en besluit tot de
        functionaalvergelijking
        \[
        \ln(ab) = \ln a + \ln b ,
        \]
        bewezen met zuivere analyse: de oppervlakte van $1$ tot $ab$ valt
        uiteen in geschaalde kopieën.
  \item Leid uit vraag 14 af: $\ln(a^n) = n\ln a$ en $\ln\frac1a = -\ln a$.
  \item \Cref{exo:g12:integ:8} leest
        $\frac{1}{n+1} + \frac{1}{n+2} + \dots + \frac{1}{2n}$ als
        rechthoeken onder $\frac{1}{1 + t}$: bereken die som voor $n = 10$
        (drie decimalen) en vergelijk met $\ln 2$. Welke sommen met een
        harmonische smaak, term voor term divergent, convergeren hier naar een
        oppervlakte?
\end{enumerate}

\medskip
\noindent\textbf{Deel IV --- Opstapelen.}
\begin{enumerate}[resume]
  \item Een auto versnelt met snelheid $v(t) = 3t^2$ m/s voor
        $t \in \intcc{0}{10}$. Bereken de afgelegde afstand, de \omterm{def:g11:stat:mean}{gemiddelde}
        snelheid, en het ogenblik waarop de ogenblikkelijke snelheid gelijk is
        aan de \omterm{def:g11:stat:mean}{gemiddelde} snelheid.
  \item Een staaf van $4$ meter heeft lineaire dichtheid
        $\rho(x) = 2 + x$ kg/m. Bereken haar totale massa en haar
        massamiddelpunt
        $\dfrac{\int_0^4 x\,\rho(x)\,\dd x}{\int_0^4 \rho(x)\,\dd x}$ --- het
        evenwichtspunt van de kartonnen driehoek uit het onderbouwvolume,
        eindelijk berekend met gewichten die variëren.
  \item De integralen van Wallis (\cref{exo:g12:integ:9}): bereken $W_0$,
        $W_1$ en $W_2 = \int_0^{\pi/2} \sin^2 t\,\dd t$ met de linearisatie
        uit \cref{pb:g12:trigo:1}. (Hun oneindige ladder klimt in de
        universitaire volumes naar een productformule voor $\pi$ en naar de
        normalisatie van de klokkromme.)
  \item Slotstuk --- de drie gezichten van het integreren: een
        \emph{oppervlakte} per definitie, een \emph{opstapeling} in de wereld
        (afstand, massa), en een \emph{\omterm{def:g12:integ:primitive}{primitieve functie}} volgens de
        hoofdstelling; en de geschiedenis ervan in drie namen. Sluit af met de
        parel van dit hoofdstuk: welke alledaagse toets op de rekenmachine is
        in het geheim de oppervlakte onder $\frac1t$ --- en welke identiteit
        bewees die oppervlakte?
\end{enumerate}
\end{problem}
